% Generated by code/tables/tabA-measurement.py. Do not edit by hand.

\begin{table}[htbp]
\centering
\begin{threeparttable}
\caption{Agreement across models and across runs}
\label{tab:agreement}
\footnotesize
\begin{spacing}{1}
\begin{tabular}{@{}lcccc@{}}
\toprule
 \\
\midrule
\multicolumn{5}{@{}l}{\emph{A. Task level: Spearman correlations of the stage-4 score across models (19,276 tasks)}} \\
 & Claude & Gemini & GPT & DeepSeek \\
Claude & 1 &  &  &  \\
Gemini & 0.900 & 1 &  &  \\
GPT & 0.865 & 0.905 & 1 &  \\
DeepSeek & 0.804 & 0.829 & 0.823 & 1 \\
Mean stage-4 score & 47.1 & 42.7 & 40.9 & 27.3 \\
\addlinespace
\multicolumn{5}{@{}l}{\emph{B. Task level: Claude Haiku 4.5 rated twice, Spearman correlation between runs on 1,000 tasks}} \\
 & Stage 1 & Stage 2 & Stage 3 & Stage 4 \\
Rating A, feasible time savings & 0.974 & 0.976 & 0.976 & 0.967 \\
Rating B, feasible completion & 0.834 & 0.860 & 0.903 & 0.931 \\
Rating C, human advantage & 0.885 & 0.958 & 0.962 & 0.947 \\
\addlinespace
\multicolumn{5}{@{}l}{\emph{C. Occupation level: Spearman correlation with the three-model score (879 occupations)}} \\
 & Gemini & GPT & DeepSeek & Four-model score \\
Stage-4 score & 0.993 & 0.987 & 0.973 & 0.999 \\
Change from stage 1 to stage 4 & 0.986 & 0.973 & 0.953 & 0.997 \\
\bottomrule
\end{tabular}
\end{spacing}
\begin{tablenotes}\notesize
\item Notes: Panel A uses the production runs on all tasks; each model's score is the composite of its own three ratings. Panel B compares the original April 2026 Claude run with a second run of the same pipeline in August 2026 on the stratified 1,000-task pilot sample. Panel C aggregates each model's task scores to occupations with the relevance weights of Section~\ref{sec:construction} and correlates them with the three-model score of the paper; the four-model score adds DeepSeek V4 Flash. The Claude column of Panel C is omitted because Claude's occupation scores enter the three-model score.
\end{tablenotes}
\end{threeparttable}
\end{table}
